% main.tex — Diplomarbeit „Decentralized Finance und die Zukunft des Finanzwesens“
% Überarbeitete und aktualisierte Fassung, 2026 — Martin Janda
%
% Kapitelinhalte werden per Pandoc aus kapitel/*.md nach build/*.tex konvertiert
% (siehe Makefile, Ziel `tex`) und hier per \input eingebunden.
% Die Kapitel-Markdown-Dateien in kapitel/ sind die alleinige Quelle für den
% Fließtext und werden NICHT in diesem Repository-Teil verändert.

\documentclass[12pt,a4paper,oneside,BCOR=0mm]{scrreprt}

% praeambel.tex — Paket- und Layout-Konfiguration
% Diplomarbeit „Decentralized Finance und die Zukunft des Finanzwesens"
% Überarbeitete Fassung 2026

% --- Encoding, Sprache, Schrift ---
\usepackage[T1]{fontenc}
\usepackage{lmodern}
\usepackage[ngerman]{babel}
\usepackage[style=german]{csquotes}
\usepackage{microtype}

% --- Zeilenabstand ---
\usepackage{setspace}
\onehalfspacing

% --- Seitenlayout ---
\usepackage[a4paper,margin=2.5cm]{geometry}

% --- Mathematik ---
\usepackage{amsmath}

% --- Tabellen ---
\usepackage{booktabs}
% Pandoc setzt Markdown-Pipe-Tabellen als longtable (braucht zusätzlich array/calc)
\usepackage{longtable}
\usepackage{array}
\usepackage{calc}

% --- Grafiken ---
\usepackage{graphicx}
\graphicspath{{abbildungen/out/}}
% Pandoc emittiert width UND height; ohne keepaspectratio würden die
% Abbildungen auf \textheight verzerrt (Pandocs eigene Preamble setzt dies auch).
\setkeys{Gin}{keepaspectratio}

% --- Bibliographie ---
\usepackage[
  backend=biber,
  style=numeric,
  sorting=nyt,
  urldate=long,
  dateabbrev=false
]{biblatex}
\addbibresource{literatur.bib}

% --- Hyperlinks (dezent, keine farbigen Rahmen) ---
\usepackage[hidelinks]{hyperref}
% xurl erlaubt Zeilenumbrüche an beliebigen Stellen in URLs (Literaturverzeichnis)
\usepackage{xurl}

% --- Pandoc-Kompatibilität ---
% \tightlist wird von Pandoc für kompakte Listen ohne Zusatzabstand erzeugt
% und ist ohne den Pandoc-Standardtemplate nicht definiert.
\providecommand{\tightlist}{%
  \setlength{\itemsep}{0pt}\setlength{\parskip}{0pt}}

% Einzelne Unicode-Zeichen aus wörtlich zitierten Roh-Extrakten (siehe
% Standard-T1-Encoding kein Glyphen-Mapping haben.
\DeclareUnicodeCharacter{03A3}{$\Sigma$}
\DeclareUnicodeCharacter{2217}{$\ast$}
% Zeichen aus dem Änderungsvermerk (Anhang B)
\DeclareUnicodeCharacter{2192}{$\rightarrow$}
\DeclareUnicodeCharacter{2194}{$\leftrightarrow$}
\DeclareUnicodeCharacter{2265}{$\geq$}

% --- Abkürzungsverzeichnis (Setup vorbereitet, Liste vorerst leer) ---
\usepackage[printonlyused]{acronym}
% Beispiel-Einträge (auskommentiert, bei Bedarf befüllen):
% \begin{acronym}
%   \acro{DeFi}{Decentralized Finance}
%   \acro{TVL}{Total Value Locked}
% \end{acronym}


\title{Decentralized Finance und die Zukunft des Finanzwesens}
\author{Martin Janda}
\date{2026}

\begin{document}

% --- Titelseite (Whitepaper-Stil; bewusst ohne Matrikelnummer/Betreuer) ---
\begin{titlepage}
  \centering
  \vspace*{2.5cm}

  {\LARGE\bfseries Decentralized Finance und die Zukunft des Finanzwesens\par}

  \vspace{0.6cm}
  {\large Eine Untersuchung möglicher Auswirkungen auf den traditionellen Finanzsektor\par}

  \vspace{1.6cm}
  {\large\scshape Whitepaper\par}

  % Titelgrafik: TVL-Verlauf 2017–2026 (Eigenanfertigung aus eingechecktem
  % DeFiLlama-Snapshot, siehe abbildungen/titelgrafik.py)
  \vspace{1.3cm}
  \includegraphics[width=0.88\textwidth]{titelgrafik.pdf}

  \vspace{1.3cm}
  {\large Martin Janda\par}
  \vspace{0.2cm}
  {\normalsize martin@janda.io\par}

  \vfill
  % Details (Einreichung/Annahme, Aktualisierungsumfang) stehen bewusst nur
  % auf der Vermerk-Seite „Über dieses Dokument“ — keine Doppelung.
  {\small Überarbeitete Fassung der Diplomarbeit, AKAD University\ · \textbf{Stand: August 2026}\par}
  \vspace*{1.2cm}
\end{titlepage}

% --- Vermerk-Seite: Herkunft, Disclaimer, Lizenz ---
\thispagestyle{empty}
\vspace*{\fill}
\noindent
\begin{minipage}{\textwidth}\small
  \textbf{Über dieses Dokument.} Dieses Whitepaper basiert auf der vom Autor
  an der AKAD University eingereichten und angenommenen Diplomarbeit
  \enquote{Decentralized Finance und die Zukunft des Finanzwesens} (Mai 2023).
  Für die Veröffentlichung wurden Marktdaten und regulatorischer Stand auf
  August 2026 aktualisiert und einzelne Kapitel überarbeitet. Prüfungsbezogene
  Bestandteile der eingereichten Fassung sind nicht Teil dieser
  Veröffentlichung.

  \medskip
  \textbf{Danksagung.} Mein besonderer Dank gilt Prof.~Dr.~Philipp Sandner~(†),
  Gründer und Leiter des Frankfurt School Blockchain Center. Als externer
  Experte im Netzwerk des FSBC habe ich in verschiedenen Projekten
  im Kryptoumfeld mit ihm zusammengearbeitet; der fachliche Austausch mit ihm hat mich
  zum Thema dieser Arbeit inspiriert. Sein früher Tod im Januar 2024 ist ein
  großer Verlust für die deutsche Blockchain-Community.

  \medskip
  \textbf{Hinweis.} Dieses Dokument dient ausschließlich Informationszwecken
  und stellt keine Anlage-, Rechts- oder Steuerberatung dar.

  \medskip
  \textbf{Lizenz.} Dieses Werk ist lizenziert unter einer Creative-Commons-Lizenz
  „Namensnennung 4.0 International“ (CC~BY~4.0),
  \url{https://creativecommons.org/licenses/by/4.0/deed.de}.

  \medskip
  \textbf{Kontakt.} Martin Janda, martin@janda.io
\end{minipage}
\vspace*{\fill}
\clearpage

% --- Vorwort zur überarbeiteten Fassung ---
\chapter*{Vorwort zur überarbeiteten Fassung}
\addcontentsline{toc}{chapter}{Vorwort zur überarbeiteten Fassung}

Diese Arbeit entstand 2023 als Diplomarbeit an der AKAD University. Kaum ein
Gegenstand verändert sich so schnell wie der hier untersuchte: Die Regulierung,
die damals erst Konturen annahm, ist mit MiCA und dem GENIUS Act Realität
geworden; die untersuchten Protokolle haben Generationswechsel, Rebrandings
und Konsolidierung durchlaufen; und mit Spot-ETFs und tokenisierten Fonds hat
die institutionelle Adoption eine neue Stufe erreicht.

Für diese Veröffentlichung habe ich die Arbeit deshalb umfassend optimiert und
auf den Stand von August 2026 gebracht: Alle Marktdaten beruhen auf einem
einheitlichen, reproduzierbaren Datenstand mit ausgewiesenen Stichtagen,
Methodik und Analyse wurden geschärft — unter anderem um ein durchgängiges
Analyseraster und eine systematische Risikobetrachtung —, und die
Entwicklungen der Jahre 2023 bis 2026 sind in neuen, gekennzeichneten
Abschnitten ergänzt, einschließlich eines Prüfstands der damaligen
Einschätzungen am tatsächlichen Verlauf.

Die Arbeit bleibt, was sie war: der Versuch, eine schnelllebige Technologie
mit wissenschaftlicher Sorgfalt einzuordnen. Wenn diese Fassung dazu beiträgt,
die Debatte um dezentrale Finanzsysteme faktenbasierter zu führen, hat sie
ihren Zweck erfüllt.

\medskip
\noindent Im August 2026\\
Martin Janda
\clearpage

% --- Verzeichnisse ---
\tableofcontents
\listoffigures
\listoftables

% Abkürzungsverzeichnis (manuell, da die Kapitel-Markdown-Quellen keine \ac-Makros
% verwenden). Enthält alle im Text verwendeten Kürzel, inkl. der laut Erstgutachten
% im Original fehlenden Einträge ETH, SPV, EH/s und wBTC.
\chapter*{Abkürzungsverzeichnis}
\addcontentsline{toc}{chapter}{Abkürzungsverzeichnis}
% longtable statt tabular: 42 Einträge passen mit 1,5-Zeilenabstand nicht mehr
% auf eine Seite — ein unteilbarer tabular erzeugte eine Leerseite davor.
\begin{longtable}{@{}p{2.6cm}p{11.8cm}@{}}
AMM   & Automated Market Maker \\
API   & Application Programming Interface (Programmierschnittstelle) \\
APY   & Annual Percentage Yield (jährliche prozentuale Rendite) \\
BaFin & Bundesanstalt für Finanzdienstleistungsaufsicht \\
BGP   & Byzantinisches Generalsproblem \\
BIS   & Bank for International Settlements (Bank für Internationalen Zahlungsausgleich) \\
BMF   & Bundesministerium der Finanzen \\
CASP  & Crypto-Asset Service Provider (Krypto-Dienstleister nach MiCA) \\
CBOT  & Chicago Board of Trade \\
CDM   & Collateralized Debt Market \\
CDP   & Collateralized Debt Position \\
CeFi  & Centralized Finance \\
CEX   & Centralized Exchange (zentralisierte Kryptobörse) \\
CFMM  & Constant Function Market Maker \\
CFTC  & Commodity Futures Trading Commission (US-Derivateaufsicht) \\
dApp  & Decentralized Application (dezentrale Anwendung) \\
DeFi  & Decentralized Finance \\
DEX   & Decentralized Exchange (dezentrale Kryptobörse) \\
EBA   & European Banking Authority (Europäische Bankenaufsichtsbehörde) \\
EH/s  & Exahashes pro Sekunde ($10^{18}$ Hashes/s) \\
ERC   & Ethereum Request for Comments (Token-Standard-Verfahren) \\
ESMA  & European Securities and Markets Authority (Europäische Wertpapieraufsicht) \\
ETF   & Exchange Traded Fund (börsengehandelter Fonds) \\
ETH   & Ether (native Kryptowährung der Ethereum-Blockchain) \\
FF    & Forschungsfrage \\
FSB   & Financial Stability Board (Finanzstabilitätsrat) \\
KMAG  & Kryptomärkteaufsichtsgesetz \\
KWG   & Kreditwesengesetz \\
KYC   & Know Your Customer (Identitätsprüfung) \\
LP    & Liquidity Provider (Liquiditätsanbieter) bzw. Liquiditätspool \\
MiCA  & Markets in Crypto-Assets (Verordnung (EU) 2023/1114) \\
NPP   & Neue-Produkte-Prozess \\
P2P   & Peer-to-Peer \\
PoS   & Proof of Stake \\
PoW   & Proof of Work \\
PSD2  & Payment Services Directive 2 (Zweite EU-Zahlungsdiensterichtlinie) \\
RWA   & Real World Assets (tokenisierte reale Vermögenswerte) \\
SEC   & U.S. Securities and Exchange Commission (US-Wertpapieraufsicht) \\
SPV   & Simplified Payment Verification \\
TradFi & Traditional Finance (traditionelles Finanzwesen) \\
TVL   & Total Value Locked \\
wBTC  & Wrapped Bitcoin \\
\end{longtable}
\clearpage

% --- Kapitel (Pandoc-generierte Fragmente aus build/, siehe Makefile) ---
\input{build/01}
\input{build/02}
\input{build/03}
\input{build/04}
\input{build/05}
\input{build/06}

% --- Anhang ---
\appendix
\input{build/90}

% --- Literaturverzeichnis ---
\printbibliography

\end{document}
